\section{收入增长：确认规模与订单语言不能互换}

2025年CNY32.557bn收入是已经在财务报表中确认的规模；新项目订单年化销售额CNY35bn以上则是公司用生命周期销售额除以生命周期年数得到的管理指标。这两个数字看似接近，却回答不同问题。前者体现已经完成交付并满足确认条件的业务，后者体现公司在当期新获项目的长期年化潜在销售规模。报告不把第二个指标加到第一个指标上，也不以其直接推导次年收入。

汽车电子项目在开发、样件、验证、定点、SOP、批量交付、客户验收和回款之间存在时间差。即使客户和产品类别都已披露，也不能在缺少车型、SOP、单价、供货份额和收入确认资料时把项目量化为某个季度的收入。本案的正确用途是：订单年化额提高对未来增长的关注度；正式财务报表的收入、分部毛利和现金流决定是否把这种关注度升级为预测信用。

\begin{exhibitbox}[复核卡：订单从何时可以进入模型]
\small
\begin{tabularx}{\textwidth}{L{2.7cm}Y Y}
\toprule
证据进展 & 可作研究结论 & 是否进入基础EPS \\
\midrule
合作/生态/产品发布 & 技术能力或商业方向 & 否 \\
定点/订单年化额 & 需求可见度和项目储备 & 否 \\
客户+车型/平台+SOP & 有条件的交付可见度 & 仍需价格/毛利资料 \\
收入/毛利/现金流披露 & 已确认经营事实 & 是，按披露期和口径进入 \\
\bottomrule
\end{tabularx}
\end{exhibitbox}

\section{智能座舱：最大的盈利池与产品升级的边界}

智能座舱2025年收入CNY20.585bn，占公司收入63.23\%，毛利率18.83\%。这使其成为本院基准盈利桥的底座。年报披露第四代座舱在理想、小米、吉利、奇瑞和广汽客户中成功配套，并列示多个国内外客户的新项目订单。该信息足以支持公司具有跨客户、跨区域交付能力，也足以说明座舱不是尚未商业化的概念。

但座舱的产品上移并未自动给出利润上移。第五代AI座舱、端云一体方案、中央计算、显示系统和HUD的产品进展具有战略价值，仍缺客户金额、车型、SOP、屏幕/域控ASP、BOM与单项目毛利。模型因此只让已披露的座舱收入和毛利影响基础盈利；新平台只在出现可审计的收入或利润字段后，才可能改变基准收入增速或估值倍数。

这种处理会使报告比市场主题叙事更保守，却避免把“未来产品架构”提前写成当前利润。对于读者，座舱的关键不是再听到多少次AI座舱的表述，而是观察分部收入、毛利、客户项目量产和营运资本是否在后续披露中共同改善。

\section{智能驾驶：增长最快的分部为何仍需要折现}

智能驾驶2025年收入CNY9.700bn，同比增32.63\%，明显快于公司总体。公司披露域控制器、传感器和算法的一体化交付能力，多个OEM实现大规模量产，并有高算力、主流和国产化平台的批量交付。对于长期竞争位置，这些事实重要：它们显示公司并非只提供单一硬件，而是能够承担平台适配、功能安全、算法、传感器和项目制造的系统工作。

同一张年报也给出限制：智能驾驶毛利率16.36\%，同比下降3.55个百分点；该分部成本增速高于收入增速。因而“智驾渗透率提高”不能成为利润率必然上升的同义词。芯片、传感器、软件开发、OEM价格谈判和项目爬坡都可能使收入弹性低于市场预期的净利弹性。

本院把智能驾驶计作已确认的增长分部，却不把它分配成某个未披露的增量净利或EPS。基准情景只使用合并收入、合并归母和稳定利润率约束；当分部毛利和现金的后续披露证明增长质量时，再决定是否提高盈利或倍数。这个顺序保证“增长”不会绕过成本和现金。

\section{网联服务与软件：有收入基础，但没有纯软件估值证据}

网联服务及其他2025年收入CNY2.272bn，说明公司确有基础软件、OTA、智能进入和服务类业务的已实现基础。它还可能增强整车项目粘性，使座舱和智驾产品从一次性交付向全生命周期服务延伸。这是系统级Tier-1相对单一硬件公司的潜在差异之一。

不过，公司没有单独披露该分部的毛利率、软件订阅/服务收费模式、续费率、客户份额或合同期限。把这部分收入直接按纯软件公司倍数估值，会忽略其可能与硬件交付、项目服务和客户协议绑定的事实。本报告因此将网联收入作为合并收入和技术粘性的证据，而不单列高倍数SOTP价值。

后续的升级证据应包括：分部毛利率、经常性服务收入占比、订阅或维护合同的收费规则、客户留存/续费及相对独立的现金流。若这些资料出现，SOTP可以从“质量校验工具”升级为更有解释力的次要估值法；在此之前，合并P/E仍是更诚实的主方法。

\section{客户链：具名客户证明什么、不证明什么}

公司年报列示了广泛的国内外客户，并对部分产品给出“成功配套”“大规模量产”“新项目订单”或“持续合作”的表述。这些措辞的强度不同。成功配套和大规模量产可以证明公司在特定产品层面已经跨过开发和交付门槛；新项目订单证明商业机会；持续合作说明关系延续。它们共同支持公司具备系统交付和客户拓展能力。

它们仍不能说明客户收入占比、每个车型的供货量、ASP、客户付款条件、项目毛利或未来续供。OEM自身销量和新车型热度也不能代替这些字段，因为一家公司可能只供应少量车型、少量功能，或经历价格变化和项目切换。报告将所有OEM列为需求锚点，而不是把品牌规模变成德赛西威收入。

客户链审计的实际价值是建立证据升级清单。每当公司、OEM或监管披露新的车型/SOP/收入资料时，研究人员可以沿同一行补充，而不是从新闻标题重新猜测。对于当前目标价，缺口意味着新增项目不会直接增加EPS，只会影响上行情景的概率。

\section{上游平台与BOM：合作关系并不等于独供关系}

Qualcomm与德赛西威的G10PH合作、NVIDIA与小鹏P7的历史项目材料，是本案中较高质量的上游平台证据。它们可以说明公司具有与全球计算平台协作、面向OEM进行系统开发的能力。公司年报同时披露国产化芯片平台已经批量交付，说明多平台路线和本地化路径并存。

但年报没有披露前五供应商名称、芯片/显示/传感器的采购价、用量、锁价、替代认证或特定平台对应的客户/车型收入。因而不应把Qualcomm或NVIDIA写为全公司确定供应商，更不能把平台合作写成独供或已实现BOM降本。上游材料CNY24.183bn、占成本91.78\%这一事实反而提示BOM是毛利敏感项。

研究上的正确反应不是忽略平台合作，也不是夸大它，而是把它放入“技术适配能力与供应链风险”两个维度。若后续出现供应中断、平台迁移或成本变化，模型应首先调整毛利和项目风险；若出现客户、SOP与收入闭环，才可讨论其带来的收入上行。

\section{营运资本与现金：汽车电子成长的第二张利润表}

2025年经营现金流CNY2.884bn，同比改善显著，支持收入并非完全停留在应计利润。与此同时，年末应收CNY9.779bn、存货CNY4.789bn，合同负债CNY1.005bn；审计报告也提示公司通常按订单生产、专用度较高，车型销售不利时可能形成库存减值风险。对于项目制汽车电子，这些余额不是背景数据，而是决定成长质量的核心变量。

2026Q1期末应收较年末下降、合同负债上升，提供了一个积极的短期观察点，但单季度仍不足以证实全年现金循环。模型不通过把库存增减或合同负债直接推导未来收入，而将它们作为中报的验证项：收入增长应与经营现金流匹配，库存和应收不应长期快于销售，减值准备不应侵蚀毛利和净利。

这也是为什么基础估值没有以很高的倍数反映远期订单。现金转化若持续改善，市场可有理由提高对增长持续性的信心；若营运资本恶化，即使收入表面增长，P/E倍数也可能收缩。

\section{全球化制造：选择权与执行成本同时存在}

公司在国内外布局研发和制造，并披露墨西哥、欧洲等海外进展。这一网络可提升服务海外客户和中国车企出海的能力，也可能缩短交付半径、改善本地供应链响应。对于全球OEM项目，制造和质量体系是进入供应链的必要条件，不能只用产品参数解释竞争位置。

海外布局也带来爬坡、合规、汇率、人员、供应链和资本开支风险。年报没有给出单厂设计产能、利用率、地区收入、毛利、现金回报或项目分配，因此不能把工厂封顶、设备安装或首次量产直接换算成未来收入。本报告仅将其列作交付韧性的潜在支撑与风险点。

升级的标准应是单厂/地区项目进入量产、收入或成本披露、利用率或现金流的方向性证据；降级的标准是延期、成本偏差或海外客户项目没有转化。这样的门槛比“海外布局加速”更可被验证。

\section{竞争格局：系统能力的优势与可替代性}

德赛西威面对的不是单一可比公司，而是全球Tier-1、国内座舱/域控供应商、OEM自研和芯片平台迁移的多重竞争。其可验证优势是系统集成、软硬件协同、功能安全、项目交付和全球制造网络；这些能力使其可以覆盖座舱、智驾和网联，而非被限制在一个零部件层级。

可替代性同样真实。OEM可能提高自研比例，其他Tier-1可能用打包报价竞争，芯片平台可能迁移，项目需求可能因车型变化被重置。没有同一车型、同一客户、同一价格的公开比较时，不能用“市场份额领先”或CR3/CR5来定量证明护城河；本案未取得口径一致的正式份额资料，因此不把份额排名当估值输入。

竞争章节的价值在于选择正确的估值语言：公司值得以已实现系统收入和项目能力被研究，但护城河必须在毛利、客户延续、现金和项目量产中持续呈现。它不是不可替代的自动结论。

\section{卖方预期分歧：EPS区间比评级数量更有用}

近90天原始报告和公开一致预期快照普遍给出正面评级，但评级数量不能说明目标价的可靠程度。更有信息量的是2026E EPS从CNY4.41到CNY5.20、2027E EPS从CNY5.12到CNY6.23的离散范围。这反映市场对收入增速、毛利、海外和产品进展的具体假设并不一致。

本院基准EPS CNY4.69处于原始报告区间的审慎一侧。它不等于对卖方预测的否定，而是把尚未闭合的客户/ASP/订单转收入和费用字段保留为折现。若正式中报使公司盈利更接近公开中位区间，基准情景的置信度提高；若低于基准，市场价格的17倍左右前瞻P/E也需要重估。

由于没有当前可审计的目标价和估值法，报告明确把Street价格锚设为零。读者可以用卖方EPS作为外部比较，却不应把过时CNY147或无来源的“共识目标价”作为当前估值支持。

\section{最终复核卡：发布后什么必须刷新}

研究报告不是永久静态文件。每次正式中报、年报、重大客户项目披露或显著价格变动后，都要重新检查价格、股本、市值、TTM/前瞻倍数、收入、毛利、现金流、营运资本、客户链和券商原始材料。刷新应遵循因果顺序：先更新事实和证据质量，再更新盈利桥，最后更新估值和行动标签。

最不应做的事情是只刷新价格和目标价而不刷新盈利前提，或只看到订单新闻就提高EPS。对德赛西威，CNY86.4基础业务价值和CNY2.49 Physical AI风险调整期权都只在本报告列示的市场时点、股本、盈利情景、概率和证据边界下成立；二者合并为CNY88.9。任何一项变动都可能让合并参考值失去可比性。

本附录的作用是提供一张复核路线图。它把研究的强结论限制在可证实的范围内，也让未来新增证据有明确的进入位置。对于高成长系统级汽车电子公司，这种克制比一次性写出更高的远期目标更具投资价值。
